\documentclass[11pt,a4paper]{article}
\usepackage[margin=25mm]{geometry}
\usepackage[T1]{fontenc}
\usepackage[utf8]{inputenc}
\usepackage{lmodern,microtype,amsmath,amssymb,amsthm,booktabs,longtable,enumitem}
\usepackage[round,sort]{natbib}
\usepackage{xurl}
\usepackage[colorlinks=true,allcolors=blue]{hyperref}
\setlength{\emergencystretch}{2em}
\numberwithin{equation}{section}
\newtheorem{theorem}{Theorem}[section]
\newtheorem{proposition}[theorem]{Proposition}
\newtheorem{lemma}[theorem]{Lemma}
\newtheorem{corollary}[theorem]{Corollary}
\newtheorem{conjecture}[theorem]{Conjecture}
\newtheorem{maintheorem}[theorem]{Theorem}
\theoremstyle{definition}
\newtheorem{definition}[theorem]{Definition}
\newtheorem{fact}{Standard fact}
\renewcommand{\thefact}{I\arabic{fact}}
\theoremstyle{remark}
\newtheorem{remark}[theorem]{Remark}
\newcommand{\Fp}{\mathbb F_p}
\newcommand{\Fq}{\mathbb F_{p^2}}
\newcommand{\Fb}{\overline{\mathbb F}_p}
\newcommand{\Z}{\mathbb Z}
\newcommand{\Q}{\mathbb Q}
\newcommand{\Res}{\operatorname{Res}}
\newcommand{\Aut}{\operatorname{Aut}}
\newcommand{\Stab}{\operatorname{Stab}}
\newcommand{\den}{\operatorname{den}_{\mathrm{red}}}
\newcommand{\leg}[2]{\left(\frac{#1}{#2}\right)}
\newcommand{\Lp}[1]{\left(\frac{#1}{p}\right)}
% Set this to the public evidence repository URL or DOI landing page before submission.
\newcommand{\DataRepositoryURL}{}

\title{Supersingular poles and signed residues\\of prime-level Fricke Hauptmoduln}
\author{David Jordan}
\date{4 November 2026}
\begin{document}
\maketitle
\begin{abstract}
For $N\in\{2,3,5,7,13\}$, let $T_N=q^{-1}+O(q)$ be the normalized Hauptmodul for
$\Gamma_0(N)^+$. For every prime $p\ge5$ with $p\nmid N$, we determine the poles and
residues of the rational function representing $T_N|U_p$ modulo $p$. Its poles are
exactly the supersingular values with effective ramification index one. If $t_i$ is
such a value, its residue multiplied by
$\prod_{k\ne i}(t_i-t_k)^{(p+1)/e_k}$ equals $-\left(\frac{-N}{p}\right)$, where the
product includes all distinct supersingular values and $e_k$ is the corresponding
ramification index. The proof combines a Kaneko--Zagier derivative argument with
the explicit evaluation
$E_{p+1}^{(p-1)/2}=\left(\frac6p\right)((p-1)/2)!\,\Delta^{(p^2-1)/24}$
on the supersingular locus. We prove this evaluation uniformly by applying an
iterated theta operator to $\eta^{p+1}$ at full level twelve. It recovers Baker's
congruence and, together with Robert's isogeny identity, gives
$\bigl(\Delta(\tau)/\Delta(\ell\tau)\bigr)^{(p^2-1)/24}=\left(\frac{\ell}{p}\right)$
at every supersingular point of $X_0(\ell)$ for every prime $\ell\ne p$.
The ramification indices also admit an interpretation as node thicknesses in the
Deligne--Rapoport Fricke quotient. Exact computations provide finite certificates
for all admissible pairs with $p\le211$.
\end{abstract}
\noindent\textbf{Keywords:} supersingular elliptic curves; Fricke groups; Hauptmoduln; Hecke operators; eta quotients; congruences.

\section{Introduction}\label{sec:intro}
The reduction of $U_p$ applied to a modular function can encode the supersingular
locus as a divisor of poles. At level one, this relation is developed by \citet{dS94} and \citet[\S10]{KZ98}; \citet{Baker98} proves the
supersingular congruence that determines the sign in the level-one residue formula.
\citet{BL18} obtain corresponding formulas for Legendre modular
polynomials. We establish a uniform signed formula on the five prime-level Fricke
curves for which $X_0(N)$ itself has genus zero. The factorial evaluation
of Theorem~\ref{thm:factorial} and the signed discriminant ratio of
Theorem~\ref{thm:C} are independent of this Fricke genus-zero setting.

The residues contain information beyond the support of the supersingular divisor.
Since the reduced function vanishes at the cusp and has only simple poles, its
residues determine the entire function. An explicit residue formula therefore
reconstructs the reduced $U_p$ image from supersingular values and their weights,
and detects vanishing when no contributing value remains. The level-one theory
provides the model for this reconstruction; the Fricke quotient introduces
additional ramification, which must enter both the pole criterion and the product.

The distinction between ramified and unramified supersingular values is essential.
A ramified value is removed from the pole divisor, but remains in the residue
product with an exponent determined by its ramification index. The algebraic proof
makes this distinction through the vanishing order of a quotient involving the Hasse
invariant and the derivative of the Hauptmodul. Section~\ref{sec:geometry} explains
the same indices as thicknesses of nodes in an integral modular model.

The level-one input is an explicit signed evaluation of
$c(E)=E_{p+1}(E)^{(p-1)/2}/\Delta(E)^{(p^2-1)/24}$ at every supersingular
curve: it equals $\left(\frac6p\right)((p-1)/2)!$. The proof applies the
Serre--Robert theta-iteration method to $F=\eta^{p+1}$ on $\Gamma(12)$.
Euler's pentagonal identity makes the quadratic twist of $F$ explicit,
and comparison with the iterated operator selects the sign without a
special CM evaluation. Constancy and Baker's squared congruence follow
from this formula.

The sign calculation also has a formulation independent of genus zero.
\citet[Theorem~1.2 and Remark~4]{Ono25} proves that the supersingular value of
$\Delta(\tau)/\Delta(\ell\tau)$ raised to $(p^2-1)/12$ is one, and exhibits
examples in which the half-exponent $(p^2-1)/24$ gives $-1$.
Theorem~\ref{thm:C} determines this remaining sign: it is
$\left(\frac{\ell}{p}\right)$ for every prime $\ell\ne p$, at every
supersingular point of $X_0(\ell)$. We deduce this discriminant-ratio
identity from the factorial evaluation and Robert's isogeny identity.
It applies without a genus-zero hypothesis, whereas the explicit residue
formula requires the coordinates fixed below.

\subsection{Notation and normalization}
Throughout the pole and residue results,
$N\in\mathcal N:=\{2,3,5,7,13\}$ and $p\ge5$ is prime with $p\ne N$. Put
\[
a_N=\frac{24}{N-1},\qquad b_N=N^{12/(N-1)},\qquad
u_N=\left(\frac{\eta(\tau)}{\eta(N\tau)}\right)^{a_N},\qquad
T_N=u_N+a_N+b_N/u_N.
\]
The pairs $(a_N,b_N)$ are $(24,4096),(12,729),(6,125),(4,49),(2,13)$.
The Fricke involution sends $u_N$ to $b_N/u_N$, and
$u_N=q^{-1}-a_N+O(q)$, so $T_N=q^{-1}+O(q)$ has zero constant term.
These are the McKay--Thompson series of classes $2A,3A,5A,7A,13A$
\citep{CN79}; the proofs below use their modular properties.

We write $T=T_N$, $u=u_N$, $q=e^{2\pi i\tau}$, $\theta=q\,d/dq$,
$h=(p-1)/2$, $k=(p^2-1)/24$, $A=\overline E_{p-1}$ and
$B=\overline E_{p+1}$.
Thus $A$ is the Hasse invariant, with $q$-expansion $1$, and $B$ has the
$q$-expansion of $E_2$ modulo $p$. Put
$\Delta=\eta^{24}$ and $\Delta_N=\Delta(\tau)\Delta(N\tau)$.
For $f=\sum c_nq^n$, define $f|U_p=\sum c_{pn}q^n$.
A bar denotes reduction modulo $p$. The finite fields are
$\Fp$, $\Fq$ and $\Fb$, and $\left(\frac{a}{p}\right)$ is the Legendre symbol.
When a modular form is evaluated at an elliptic curve, an invariant differential
is understood; all weight-zero ratios are independent of its choice.

The functions $u$ and $T$ remain coordinates on $X_0(N)_{\Fp}$ and
$X_0(N)^+_{\Fp}$ respectively. Indeed, the eta quotient is a modular unit with
divisor the difference of the two cusp sections, with unit leading coefficients
at both cusps away from $N$. It therefore gives the integral genus-zero coordinate;
the quotient by $u\mapsto b_N/u$ has coordinate $T$ and a single cusp.
This also fixes the normalization used in every product below.
\subsection{Weights}
Let $t_1,\dots,t_s\in\Fb$ be the distinct supersingular values of $T$. For a point
$x=(E,C)$ of $X_0(N)_{\Fb}$ over $t_i$ put
\[
m_x=\bigl|\Aut(E,C)/\{\pm1\}\bigr|,\qquad e_i=m_x\,\bigl|\Stab_{\langle w_N\rangle}(x)\bigr|.
\]
Call $t_i$ \emph{thin} if $e_i=1$ and put $D_{\rm thin}(T)=\prod_{e_i=1}(T-t_i)$. For thin $t_i$ let
\[
r_i=\Res_{T=t_i}R_{N,p}(T),\qquad W_i=\prod_{k\ne i}(t_i-t_k)^{(p+1)/e_k}.
\]
Section~\ref{sec:weights} shows that $e_i$ does not depend on the representative of the Fricke
orbit, that it is the effective (tame) inertia of the quotient map at $x$ (not the order of a group
containing $\pm1$), and that every exponent $(p+1)/e_k$ is an integer.

\subsection{Main results}
\begin{maintheorem}\label{thm:main}
Let $N\in\{2,3,5,7,13\}$ and let $p\ge5$ be a prime with $p\ne N$.
\begin{enumerate}[label=\textup{(\alph*)}]
\item \textbf{(Rationality.)} There is a unique $R_{N,p}\in\Fp(T)$ with $\overline{T|U_p}=R_{N,p}(\overline T)$. It vanishes at $T=\infty$.
\item \textbf{(Pole support.)} $R_{N,p}=P/D_{\rm thin}$ with $P\in\Fp[T]$,
$\deg P<\deg D_{\rm thin}$ and $\gcd(P,D_{\rm thin})=1$. Every pole is simple and has a nonzero residue. In particular $\den R_{N,p}=D_{\rm thin}$.
\item \textbf{(Residue as a Hasse derivative.)} If $t_i$ is thin and $x$ lies over $t_i$, then
\[
r_i=\frac{144\,(\theta T)(x)^{p+1}}{B(x)^2},\qquad \theta=q\frac{d}{dq},
\]
a ratio of two mod-$p$ modular forms of weight $2(p+1)$, evaluated at $x$.
\item \textbf{(Product.)} For thin $t_i$, with $x=(E,C)$ over $t_i$ and $E'=E/C$,
\[
r_iW_i=\frac{B^{p-1}}{\Delta_N^{\,k}}(x)=\Lp{N}\,c(E)\,c(E'),\qquad
c(E):=\frac{B(E,\omega)^{(p-1)/2}}{\Delta(E,\omega)^{k}}.
\]
\item \textbf{(Signed weighted residue formula.)} For every thin $t_i$,
\[
\boxed{\;r_iW_i=-\Lp{-N}\;}\qquad\text{equivalently}\qquad
r_i=-\Lp{-N}\prod_{k\ne i}(t_i-t_k)^{-(p+1)/e_k}.
\]
In particular $r_i\in\Fp^\times$ and $W_i\in\Fp^\times$, also when $t_i\notin\Fp$.
\item \textbf{(Hecke form.)} With $\Psi_{N,p}(X,Y)=(X^p-Y)(X-Y^p)+pH_{N,p}(X,Y)$ as in
Proposition~\ref{prop:Psi}, $\overline H_{N,p}(t^p,t)=\Lp{-N}W_i^{-1}$ if $t=t_i$ is thin, and
$\overline H_{N,p}(t^p,t)=0$ for every other $t\in\Fq$.
\end{enumerate}
\end{maintheorem}

\begin{corollary}[reconstruction]\label{cor:reconstruction}
The supersingular values and their indices determine the entire reduced function:
\begin{equation}\label{eq:partial-fractions}
R_{N,p}(T)=-\Lp{-N}\sum_{e_i=1}\frac{1}{W_i(T-t_i)}.
\end{equation}
Equivalently, if $D_{\rm thin}=\prod_f f$ is its factorization into monic
irreducibles over $\Fp$, then
\begin{equation}\label{eq:factor-reconstruction}
R_{N,p}(T)=\sum_{f\mid D_{\rm thin}} c_f\frac{f'(T)}{f(T)},\qquad
c_f=-\Lp{-N}W_i^{-1}\in\Fp^\times\quad\text{for any root }t_i\text{ of }f.
\end{equation}
Both sums are empty when $D_{\rm thin}=1$.
\end{corollary}
\begin{proof}
Subtracting the principal parts leaves a polynomial, which is zero because the
function and each summand vanish at infinity. Frobenius-conjugate poles have
conjugate residues; the residues lie in $\Fp$ by Theorem~\ref{thm:main}(e), so
they are equal on each irreducible factor. The identity
$f'/f=\sum_{f(t_i)=0}(T-t_i)^{-1}$ gives the second formula.
\end{proof}

\begin{theorem}[supersingular factorial evaluation]\label{thm:factorial}
For every prime $p\ge5$, every supersingular elliptic curve $E/\Fb$ and
every nonzero invariant differential $\omega$,
\begin{equation}\label{eq:factorial-evaluation}
B(E,\omega)^h=\Lp6\,h!\,\Delta(E,\omega)^k.
\end{equation}
Equivalently, in the graded ring $\Fp[E_4,E_6]$,
\begin{equation}\label{eq:factorial-divisibility}
A\mid B^h-\Lp6\,h!\,\Delta^k.
\end{equation}
\end{theorem}

\begin{corollary}[constancy and Baker's congruence]\label{cor:D}
The value $c(E)$ is the same for all supersingular $E/\Fb$, namely
$c_p=\Lp6\,h!\in\Fp^\times$. In particular
$c_p^2=-\Lp{-1}$, and
\[
E_{p+1}^{p-1}\equiv-\Lp{-1}\Delta^{(p^2-1)/12}
\pmod{(p,E_{p-1})}.
\]
\end{corollary}

For a prime $\ell$, let $Y_0(\ell)$ denote the non-cuspidal modular
curve over $\Z[1/\ell]$. For $x=(E,C)$, let $\pi:E\to E'=E/C$ and
choose a nonzero invariant differential $\omega'$ on $E'$. Define
\begin{equation}\label{eq:Dmoduli}
D_\ell(E,C):=\ell^{12}
\frac{\Delta(E,\pi^*\omega')}{\Delta(E',\omega')}.
\end{equation}
The differential weights make this independent of $\omega'$ and of the
representative of $x$. It is a unit on $Y_0(\ell)$ over $\Z[1/\ell]$,
with $q$-expansion $\Delta(q)/\Delta(q^\ell)$: on the Tate curve with
$C=\mu_\ell$, $\pi^*\omega'_{\rm can}=\ell\omega_{\rm can}$.
Its complex realization is therefore $\Delta(\tau)/\Delta(\ell\tau)$.
We use the same notation for its reduction at $p\ne\ell$.

\begin{theorem}[signed discriminant ratio]\label{thm:C}
Let $p\ge5$ and $\ell\ne p$ be primes. At every supersingular point
$x$ of $Y_0(\ell)_{\Fb}$,
\[
D_\ell(x)^{(p^2-1)/24}=\leg{\ell}{p}.
\]
For $\ell=N\in\mathcal N$, $D_N=u_N^{N-1}$, so this reads
$u_N(x)^{(p^2-1)/a_N}=\Lp{N}$.
\end{theorem}

Section~\ref{sec:factorial} proves the factorial evaluation and
Corollary~\ref{cor:D}. Section~\ref{sec:proof} deduces
Theorem~\ref{thm:C} and completes the signed residue formula.

\subsection{Consequences}\label{subsec:consequences}
The reconstruction in Corollary~\ref{cor:reconstruction} gives concrete formulas
once the supersingular values and their effective indices are known.
The connection to Nakaya's squarefree Fricke supersingular polynomials is a
change of coordinate: with the normalizations in \citet[\S4]{Nakaya19},
\begin{equation}\label{eq:nakaya-coordinate}
\mathcal S_{N,p}(T):=\prod_i(T-t_i)=ss_p^{(N*)}(T+c_N),\qquad
\begin{array}{c|rrrrr}
N&2&3&5&7&13\\\hline
c_N&104&42&16&9&0
\end{array}
\end{equation}
For $N=2,3,5,7$, Nakaya's printed eta formulas have additive
constants $128,54,22,13$, giving the shifts displayed above. At level
$13$ his convention is zero constant term in the $q$-expansion; this
identifies the normalized Hauptmodul with $T_{13}$ and gives $c_{13}=0$.
All constants are reduced modulo $p$; $\mathcal S_{N,p}$ is distinct from the
weighted polynomial $S$ of Proposition~\ref{prop:S}. Thus
$D_{\rm thin}=\mathcal S_{N,p}/\prod_{e_i>1}(T-t_i)$, while the omitted
values still occur in the weights $W_i$.
The exceptional values can also be specified by a polynomial independent of $p$.
Let $h_N$ be the elliptic-value polynomial in Table~\ref{tab:elliptic-values}.
Lemma~\ref{lem:elliptic-specialization} and Theorem~\ref{thm:main}(b) give
\begin{equation}\label{eq:elliptic-gcd}
D_{\rm thin}(T)=\frac{\mathcal S_{N,p}(T)}
{\gcd\bigl(\mathcal S_{N,p}(T),\overline h_N(T)\bigr)},
\end{equation}
where the gcd is monic. Thus the denominator is obtained from Nakaya's
supersingular polynomial by removing its common factors with the reduction
of a fixed characteristic-zero elliptic-value polynomial.

At level two, put
$\varepsilon_p=(1-\left(\frac{-1}{p}\right))/2$ and
$\nu_p=(1-\left(\frac{-2}{p}\right))/2$.
The level-two indices described in Section~\ref{sec:geometry} give
\begin{equation}\label{eq:leveltwo-denominator}
D_{2,p}(T)=\frac{ss_p^{(2*)}(T+104)}
{(T+104)^{\varepsilon_p}(T-152)^{\nu_p}}.
\end{equation}
The two removed values correspond to $0$ and $256$ in Nakaya's coordinate;
they are distinct for $p\ge5$.

\begin{corollary}[level-two degree and vanishing]\label{cor:leveltwo-vanishing}
For every prime $p\ge5$,
\begin{equation}\label{eq:leveltwo-degree}
\deg D_{2,p}=\frac{p-4+\left(\frac{-1}{p}\right)+2\left(\frac{-2}{p}\right)}8.
\end{equation}
In particular $T_{2A}|U_p\equiv0\pmod p$ if and only if $p=5$ or $7$.
\end{corollary}
\begin{proof}
The two possible thick values have indices four and two, and supersingularity
indicators $\varepsilon_p$ and $\nu_p$, respectively. Corollary~\ref{cor:degree}
therefore gives
$\deg D_{2,p}=(p-1)/8-\varepsilon_p/4-\nu_p/2$, which is
\eqref{eq:leveltwo-degree}. This degree is zero at $p=5,7$ and is at least
$(p-7)/8>0$ for $p\ge11$. Theorem~\ref{thm:main}(b) proves the last assertion.
\end{proof}

\citet[Lemma~4.8 and Table~4.3]{CMT19} establish $p$-adic
annihilation of $T_{2A}$ at $p=5,7$. The corollary specifies the exact
first-step modulo-$p$ vanishing criterion. Table~\ref{tab:leveltwo-examples}
illustrates the nonzero cases; at $p=17$, for example,
$R_{2,17}=D_{2,17}'/D_{2,17}$ and both residues are one.

\begin{table}[htbp]\centering
\begin{tabular}{rll}\toprule
$p$ & $D_{2,p}(T)$ & $P_{2,p}(T)$\\\midrule
11 & $T-5$ & $1$\\
13 & $T-3$ & $-2$\\
17 & $T^2-T-8$ & $2T-1$\\
29 & $T^3-12T^2+12T-6$ & $4T^2-5T-8$\\\bottomrule
\end{tabular}
\caption{Level-two examples, with $R_{2,p}=P_{2,p}/D_{2,p}$.
Coefficients use the unique integer representatives in $(-p/2,p/2)$.
These entries are included in the certificates of Section~\ref{sec:comp}.}
\label{tab:leveltwo-examples}
\end{table}

\subsection{Relation to earlier work}
The derivative argument in Section~\ref{sec:kz} follows the level-one method
presented by \citet[\S10]{KZ98}, who trace it to Koike's work
on Atkin's congruences \citep{Koike73}. The period formula of \citet{dS94}
and the congruence of \citet{Baker98} supply the level-one precedents for the
residue product and its sign; \citet{BL18} treat the corresponding Legendre
modular polynomials. The isogeny transformation of $B$ is due to
\citet[Th\'eor\`eme B, p.~106; \S3.3, Lemma~7 and (19), p.~124]{Robert80};
we include a $q$-expansion proof to fix the differential normalization.
Here the level-one factor $-\left(\frac{-1}{p}\right)$ is recovered from
Theorem~\ref{thm:factorial} and Wilson's theorem.

The factorial evaluation uses a classical mechanism: Serre's functional
equation for Eisenstein series, presented by
\citet[\S2.5, Th\'eor\`eme~2, p.~119]{Robert80}, produces factorial
coefficients by iterating theta and then evaluating on the supersingular
locus. We apply that mechanism to the form $\eta^{p+1}$ of weight
$(p+1)/2$ at full level twelve. Its explicit quadratic twist gives the
signed value of $c_p$ simultaneously in all residue classes of $p$
modulo twelve. Robert's cited theorem concerns pairs of even-weight
Eisenstein series with complementary weights summing to $p+1$; it does
not state the factorial evaluation of Theorem~\ref{thm:factorial}. The argument here does not use
Baker's congruence or the signed discriminant-ratio identity as an input.

Fricke supersingular polynomials are studied by \citet{Sakai11,Sakai14},
\citet{Nakaya19} and \citet{Morton23}. These papers treat
orthogonal-polynomial representations, factorization and resultant
identities. Formula~\eqref{eq:nakaya-coordinate} identifies the supersingular
values used here with that literature; the present residue formula additionally
assigns a nonzero coefficient to each thin pole and retains the ramified values
in its weighted product.

\citet[\S4.2, Table~1 and Appendix~C]{JST19} determine the elliptic values
of normalized Fricke Hauptmoduln and their associated class fields. Their
$j_N$ equals our $T_N$. Table~\ref{tab:elliptic-values} records their
polynomials at the five levels considered here. Their Schwarzian equation
detects the same local ramification indices that enter our Hasse-derivative
calculation; \eqref{eq:elliptic-gcd} expresses its elliptic-value locus as
the locus of possible supersingular pole cancellations. The signed residue
formula concerns the remaining thin values.

\citet[Proposition~3.1 and \S4]{DS26} use Deligne's supersingular
expansion of $pJ|U_p$ and prime-power level modular functions to
characterize the monstrous exponents. Here we determine the signed
residues on five fixed Fricke curves.

\citet[Theorem~1.1]{BGKMR26} study iterated $U_p$ congruences for
positive-weight meromorphic cusp forms on $\mathrm{SL}_2(\Z)$ whose
interior pole reduces to a supersingular point. Their $p$-adic
approximation by holomorphic cusp forms concerns a different setting
from our first-step modulo-$p$ pole divisor and signed residue formula.

\citet[Theorem~3, p.~72; \S4]{Landweber88} gives the
half-exponent identity
$\Delta^{(p^2-1)/24}=(e_a-e_b)^{(p^2-1)/4}$ for supersingular
models $y^2=4\prod_{a=1}^3(x-e_a)$ and $a\ne b$.
Theorem~\ref{thm:C} instead compares discriminants along a prime-degree
isogeny, with its differential normalization explicit, and determines
the sign from Theorem~\ref{thm:factorial} and Robert's identity.
\citet[Theorem~1.2 and Remark~4]{Ono25}
proves the squared isogeny-ratio identity and records examples with nontrivial
half-exponent sign. His work develops the supersingular congruences arising in
the description of the divided beta family by \citet[Theorem~1.2]{Behrens09}
and the explicit calculations of \citet[\S5.4, Conjecture~5.17]{Larson20}.
These connections motivate the discriminant identity; no new statement about the
Adams--Novikov spectral sequence is needed here.

Our main result combines the Hasse-derivative calculation with the explicit
factorial evaluation to determine the signed residue products on the five
Fricke curves. The latter evaluation also gives the discriminant-ratio
identity for every prime isogeny degree different from $p$.
The exponent $(p^2-1)/24$ and the isogeny transformation of $E_{p+1}$
are classical inputs; the formulas above specify their signed consequences.

\section{Standard inputs and the residue criterion}\label{sec:imports}
We use the $q$-expansion principle and Kodaira--Spencer isomorphism for prime-to-$p$
level structures \citep{Katz73,KM85,Gross90}. On a fine moduli cover, the Hasse invariant
has simple zeros precisely at the supersingular points, and Kodaira--Spencer identifies
a weight-two form $f$ with $f\,dq/q$. In particular, $\theta T$ represents $dT$.
We record the statements used repeatedly.
\begin{fact}\label{fact:modp}
\begin{enumerate}[label=\textup{(\roman*)}]
\item The $q$-expansion principle holds at prime-to-$p$ level. The forms
$A$ and $B$ have $q$-expansions $1$ and $E_2$ modulo $p$, respectively.
\item The Hasse invariant $A$ has simple zeros precisely at the
supersingular points on a fine moduli cover.
\item For $p\ge5$ the stabilizers relevant to $X_0(N)$ and its Fricke
quotient are tame, and these coarse quotients commute with reduction.
The supersingular values of $T$ lie in $\Fq$.
\item Kodaira--Spencer identifies a weight-two form $f$ with
$f\,dq/q$; in particular $\theta T$ represents $dT$.
\end{enumerate}
\end{fact}
The field-of-definition assertion in (iii) follows by choosing a model
with $p^2$-Frobenius $[-p]$ as follows; every cyclic subgroup of order
$N$ is then rational as a subgroup scheme.

\begin{lemma}[Deuring models]\label{lem:deuring}
Let $E_1/\Fb$ be supersingular, $p\ge5$. Then $E_1$ has a model over $\Fq$ whose $p^2$-power
Frobenius endomorphism is $[-p]$. For such a model every finite subgroup of order prime to $p$ is
$\Fq$-rational, and $E_1/K$ again has Frobenius $[-p]$.
\end{lemma}
\begin{proof}
A supersingular curve admits a model over $\Fq$
\citep{Deuring41,Waterhouse69}. For any such model $E/\Fq$, the $p$-power Frobenius $F:E\to E^{(p)}$ and the
Verschiebung $V:E^{(p)}\to E$ are both purely inseparable of degree $p$ (the Hasse invariant
vanishes), so $F^{(p)}:E^{(p)}\to E^{(p^2)}=E$ and $V$ differ by an automorphism $u$ of $E$:
$F^{(p)}=u\circ V$, whence $\pi_{p^2}=F^{(p)}\circ F=u\circ V\circ F=u\circ[p]$. Thus
$\pi_{p^2}=p\,u$ with $u\in\Aut(E)$. If $j\ne0,1728$ then $u=\pm1$ and the quadratic twist changes
the sign; if $j=0$ (resp.\ $1728$) the curve $y^2=x^3+1$ (resp.\ $y^2=x^3+x$) over $\Fp$ has
trace $0$, so $\pi_{p^2}=\pi_p^2=[-p]$ \citep{Deuring41,Waterhouse69}; see also
\citet[Lemma~4.3]{Ono25}. Since $[-p]$ is central, Galois preserves every subgroup, and the quotient
isogeny $\phi$ satisfies $\pi'\phi=\phi\pi=\phi[-p]=[-p]\phi$, so $\pi'=[-p]$.
\end{proof}

\begin{proposition}\label{prop:Psi}
Let $\Delta_p$ be the set of integral matrices of determinant $p$ with lower-left entry divisible by
$N$. There is a unique $\Psi_{N,p}\in\Z[X,Y]$ with
\[
\Psi_{N,p}(X,T(\tau))=(X-T(p\tau))\prod_{b=0}^{p-1}\Bigl(X-T\bigl(\tfrac{\tau+b}{p}\bigr)\Bigr).
\]
It is monic of degree $p+1$ in each variable, symmetric, and satisfies
$\Psi_{N,p}\equiv(X^p-Y)(X-Y^p)\pmod p$.
\end{proposition}
\begin{proof}
The matrices $\mathrm{diag}(p,1)$ and $\begin{pmatrix}1&b\\0&p\end{pmatrix}$, $0\le b<p$, represent
$\Gamma_0(N)\backslash\Delta_p$. Conjugation by $W_N=\begin{pmatrix}0&-1\\N&0\end{pmatrix}$ maps
$\begin{pmatrix}a&b\\c&d\end{pmatrix}$ to $\begin{pmatrix}d&-c/N\\-Nb&a\end{pmatrix}$, so it preserves
$\Delta_p$. Hence the elementary symmetric functions of the roots are $\Gamma_0(N)^+$-invariant,
holomorphic on the upper half-plane, with a pole only at the single cusp.
Their coefficients initially lie in $\Z[\zeta_p]((q^{1/p}))$.
Translation of $\tau$ removes fractional powers, and cyclotomic Galois automorphisms
permute the roots, so the expansions lie in $\Z((q))$.
Subtracting integral powers of $T=q^{-1}+O(q)$ then puts them in $\Z[T]$. Modulo $1-\zeta_p$ all $T((\tau+b)/p)$ coincide with $T(q^{1/p})$, whose $p$-th power is $T(q)$, and
$T(p\tau)\equiv T^p$; this gives the congruence. For irreducibility, the reduction of $\Gamma_0(N)$ modulo $p$ is
$\mathrm{SL}_2(\Fp)$, by the Chinese remainder theorem and the surjectivity of reduction
for $\mathrm{SL}_2(\Z)$. It acts transitively on the $p+1$ cyclic subgroups of order $p$.
The $p+1$ roots displayed above are distinct: their leading terms are $q^{-p}$ and
$\zeta_p^{-b}q^{-1/p}$ for $0\le b<p$. They therefore form one orbit, so
$\Psi_{N,p}(X,T)$ is irreducible over $\mathbb C(T)$.
The pole order of the $j$th elementary symmetric function is at most
$p+(j-1)/p\le p+1$; the constant coefficient has pole order $p+1$.
Thus the polynomial has degree $p+1$ in $Y$ as well as in $X$.
For symmetry, $\mathrm{diag}(1,p)\in\Delta_p$ gives
$\Psi_{N,p}(T(\tau),T(p\tau))=0$ by evaluating the defining relation at $p\tau$.
Irreducibility and Gauss's lemma therefore identify the reversed integral primitive
polynomial with the original up to sign.
Hence $\Psi_{N,p}(Y,X)=c\,\Psi_{N,p}(X,Y)$ with $c=\pm1$. Reducing modulo $p$, where both polynomials
equal the symmetric $(X^p-Y)(X-Y^p)$, forces $c=1$.
\end{proof}

Define $H=H_{N,p}\in\Z[X,Y]$ by $\Psi_{N,p}=(X^p-Y)(X-Y^p)+pH$; $H$ is symmetric. Put
\[
\psi=\frac{T(q)^p-T(q^p)}{p}\in\Z((q)),
\]
which is integral because $T(q^p)\equiv T(q)^p\pmod p$.

\begin{proposition}[residue criterion]\label{prop:res}
In $\Fp((q))$,
\[
\overline\psi=\frac{\overline H(T^p,T)}{T^{p^2}-T},\qquad
\overline{T|U_p}=\overline Q(T)+\overline\psi,
\]
where $\overline Q\in\Fp[T]$ is defined by $-[X^p]\Psi_{N,p}=Y^p+pQ(Y)$. Hence
$R_{N,p}=\overline Q(T)+\overline H(T^p,T)/(T^{p^2}-T)$ exists and is unique. It vanishes at
$\infty$, its poles are simple and lie in $\Fq$, and for $t\in\Fq$
\[
\Res_{T=t}R_{N,p}=-\overline H(t^p,t).
\]
\end{proposition}
\begin{proof}
Substitute $X=T(q^p)=T^p-p\psi$ and $Y=T$ into
$\Psi_{N,p}=0$. Since $X-Y^p=-p\psi$, division by $p$ and reduction give
$-\overline\psi\,(T^{p^2}-T)+\overline H(T^p,T)=0$. The sum of the roots of $\Psi_{N,p}(X,T)$ is
$T(q^p)+p\,(T|U_p)=T^p+pQ(T)$, so $T|U_p=Q(T)+\psi$. The polynomial $T^{p^2}-T$ is squarefree, its
roots are exactly $\Fq$, and its derivative is $-1$. $R_{N,p}$ vanishes at $\infty$ because
$T|U_p\in q\Z[[q]]$ (the constant term of $T$ is $0$). Uniqueness holds because $\overline T$ is
transcendental over $\Fp$.
\end{proof}

\begin{corollary}\label{cor:rFp}
For every $t\in\Fq$, $\Res_{t}R_{N,p}\in\Fp$ and $\Res_{t^p}R_{N,p}=\Res_tR_{N,p}$.
\end{corollary}
\begin{proof}
$(\Res_tR)^p=\Res_{t^p}R=-\overline H(t^{p^2},t^p)=-\overline H(t,t^p)=-\overline H(t^p,t)=\Res_tR$, using
$t^{p^2}=t$, the $\Fp$-rationality of $R$ and the symmetry of $H$.
\end{proof}

Proposition~\ref{prop:res} gives Theorem~\ref{thm:main}(a) and the a~priori bound
$\den R_{N,p}\mid T^{p^2}-T$ used by the certificates in Section~\ref{sec:comp}.

\section{Weights and the weighted supersingular polynomial}\label{sec:weights}

\begin{lemma}[effective inertia]\label{lem:inertia}
Let $x=(E,C)$ be a non-cuspidal point of $X_0(N)_{\Fb}$, supersingular or ordinary. Let $G_x$ be the group of automorphisms
of $(E,C)$ together with the isomorphisms $(E,C)\to w_N(E,C)=(E/C,E[N]/C)$ when $w_N(x)=x$. It is a group once
$w_N^2(E,C)\cong(E,C)$ is identified via $[N]$.
\begin{enumerate}[label=\textup{(\roman*)}]
\item $|G_x/\{\pm1\}|=m_x|\Stab_{\langle w_N\rangle}(x)|=e_x$. The value is the same at $x$ and at $w_N(x)$.
\item $\{\pm1\}$ is exactly the kernel of the action of $G_x$ on the deformation space of $x$. Since
$e_x\in\{1,2,3,4,6\}$ is prime to $p$, $G_x/\{\pm1\}$ is cyclic and acts faithfully on the tangent line.
Hence, for a local parameter $z$ of a fine moduli cover at a point over $x$, $T-t=z^{e_x}\cdot(\text{unit})$.
\item If $x$ is supersingular, $(p+1)/e_x\in\Z$.
\end{enumerate}
\end{lemma}
\begin{proof}
(i) If $w_N(x)=x$, the isomorphisms $(E,C)\to w_N(E,C)$ form a torsor under $\Aut(E,C)$, so
$|G_x|=2|\Aut(E,C)|$. Conjugation by the isogeny $E\to E/C$ identifies $\Aut(E,C)$ with
$\Aut(E/C,E[N]/C)$, so $m_x=m_{w_N(x)}$.
(ii) $[-1]$ acts trivially on deformations. A Fricke-type element acts
on the deformation functor by
$(\widetilde E,\widetilde C)\mapsto
(\widetilde E/\widetilde C,\widetilde E[N]/\widetilde C)$,
followed by the chosen identification of $w_N(E,C)$ with $(E,C)$.
If its action were trivial, the universal deformation would admit an
isomorphism $\widetilde E/\widetilde C\cong\widetilde E$.
Composing with the quotient map would give a degree-$N$ endomorphism
of its generic fibre, whose endomorphism ring is $\Z$; this is impossible
since $N$ is not a square. Similarly, a nontrivial automorphism other
than $\pm1$ acting trivially would extend to the universal deformation
and contradict that generic endomorphism ring. The same argument works at ordinary points. A finite group of order prime to $p$ acting faithfully on $k[[z]]$ is cyclic and
acts faithfully on $(z)/(z^2)$, so the invariant ring is $k[[z^{e}]]$.
(iii) $e=3$ or $6$ forces $m_x=3$, so $j=0$ is supersingular and $p\equiv2\pmod3$; $e=4$ forces
$m_x=2$, so $j=1728$ is supersingular and $p\equiv3\pmod4$. For $p\ge5$, $m_x\in\{1,2,3\}$.
\end{proof}

For explicit calculations, $m_x$ can be read off from the multiplicity $\mu$ of $u(x)$ as a root of
$P_N(u)-j\,u^N$: $m_x=3/\mu$ over $j=0$, $m_x=2/\mu$ over $j=1728$, and $m_x=1$ (with $\mu=1$) otherwise.
These are ramification indices of the coarse map $X_0(N)\to X(1)$.
The Fricke stabilizer has order two exactly when $u(x)^2=b_N$.

\begin{lemma}[elliptic values and specialization]\label{lem:elliptic-specialization}
Let $h_N(T)=\prod_{\tau}(T-T_N(\tau))$, where $\tau$ runs over the
inequivalent elliptic points of $\Gamma_0(N)^+$. For the five levels under
consideration, $h_N$ is given in Table~\ref{tab:elliptic-values}.
For every prime $p\ge5$ with $p\ne N$, its reduction is squarefree, and
its roots in $\Fb$ are exactly the values with effective index $e>1$.
Consequently,
\[
\gcd(\mathcal S_{N,p},\overline h_N)=\prod_{e_i>1}(T-t_i).
\]
\end{lemma}
\begin{proof}
The characteristic-zero values and elliptic orders are those of
\citet[\S4.2, Table~1 and Appendix~C]{JST19}, in the same normalization
$T_N=q^{-1}+O(q)$. Their elliptic orders are the effective indices of
Lemma~\ref{lem:inertia} in characteristic zero.
We justify the specialization of their support. A nontrivial effective
stabilizer arises either from an automorphism of $(E,C)$ or from a Fricke
fixed point. The former can occur only over $j=0$ or $1728$. Since
$p\ge5$ and $p\nmid N$, the corresponding elliptic curve and its
automorphisms lift to characteristic zero, and the prime-to-$p$ cyclic
subgroup $C$ lifts uniquely over a strict henselian base. An automorphism
preserving $C$ on the special fibre therefore preserves its lift.
Conversely, these prime-to-$p$ automorphisms persist under reduction.
Fricke fixed points are given by $u^2=b_N$, which is finite \'etale away
from $2N$ and has the same fixed-point locus under specialization.
This accounts for all values with $e>1$ in either characteristic.
Directly from the displayed polynomials, the discriminants of $h_N$, in
the order $N=2,3,5,7,13$, are
\[
2^{16},\qquad 2^4\cdot3^6,\qquad 2^{12}\cdot5^3,\qquad
2^4\cdot3^6\cdot7^2,\qquad 2^{12}\cdot3^6\cdot13.
\]
Thus distinct elliptic values do not collide at any admissible prime.
The roots of $\overline h_N$ are precisely the values with $e>1$;
intersecting this set with the supersingular locus proves the gcd formula.
\end{proof}

\begin{table}[htbp]\centering
\begin{tabular}{cll}\toprule
$N$ & $h_N(T)$ & Effective indices\\\midrule
2 & $(T+104)(T-152)$ & $4,\ 2$\\
3 & $(T+42)(T-66)$ & $6,\ 2$\\
5 & $(T+16)(T^2-12T-464)$ & $2,\ 2$\\
7 & $(T+10)(T+9)(T-18)$ & $2,\ 3,\ 2$\\
13 & $(T+4)(T+3)(T^2-4T-48)$ & $2,\ 3,\ 2$\\\bottomrule
\end{tabular}
\caption{Elliptic-value polynomials from \citet{JST19}.
The indices are listed in the order of the displayed factors; an index
beside a quadratic factor applies to both roots. At admissible primes,
the supersingular roots of $\overline h_N$ are exactly the thick values
removed in \eqref{eq:elliptic-gcd}.}
\label{tab:elliptic-values}
\end{table}

\begin{lemma}\label{lem:AW}
As mod-$p$ forms on $\Gamma_0(N)$: $A|_{p-1}W_N=\Lp{N}A$, $\Delta_N|_{24}W_N=\Delta_N$ and
$\theta T|_2W_N=\theta T$.
\end{lemma}
\begin{proof}
For a level-one form $f$ of weight $k$, $f|_kW_N=N^{k/2}f(N\tau)$. The form $E_{p-1}(N\tau)$ has
$q$-expansion $\equiv1\equiv E_{p-1}(\tau)$, so it equals $A$ modulo $p$ by Fact~\ref{fact:modp}(i);
and $N^{(p-1)/2}\equiv\Lp N$. The other two identities hold in characteristic $0$:
$\Delta(-1/(N\tau))\Delta(-1/\tau)=N^{12}\tau^{24}\Delta(N\tau)\Delta(\tau)$, and $T$ is
$w_N$-invariant.
\end{proof}

\begin{proposition}[weighted supersingular polynomial]\label{prop:S}
Let $S(T)=\prod_{k=1}^s(T-t_k)^{(p+1)/e_k}$. Then $S\in\Fp[T]$ and
\[
S(\overline T)=\Delta_N^{-k}=q^{-(N+1)k}\prod_{n\ge1}(1-q^n)^{-24k}(1-q^{Nn})^{-24k}\quad\text{in }\Fp((q)).
\]
Equivalently, $S=A^{p+1}/\Delta_N^{k}$ as functions on $X_0(N)^+_{\Fp}$. In particular
$\sum_k e_k^{-1}=(N+1)(p-1)/24$.
\end{proposition}
\begin{proof}
$F:=A^{p+1}/\Delta_N^k$ is a ratio of mod-$p$ forms of weight $p^2-1$. By Lemma~\ref{lem:AW} it is
$w_N$-invariant, since $\Lp N^{p+1}=1$. $\Delta_N$ does not vanish on $Y_0(N)$, so $F$ is regular off the cusp of
$X_0(N)^+$. Its $q$-expansion there is $\Delta_N^{-k}=q^{-(N+1)k}(1+O(q))$, because $A$ has
$q$-expansion $1$. Hence $F$ is a monic polynomial in $T$ of degree $(N+1)k$. By
Fact~\ref{fact:modp}(ii) and Lemma~\ref{lem:inertia}(ii), $F$ vanishes exactly at the supersingular
$t_k$, to order $(p+1)/e_k$ in $T-t_k$. Comparing degrees gives the mass identity.
\end{proof}

\begin{corollary}\label{cor:degree}
The degree of the reduced denominator is
\[
\deg D_{\rm thin}=\frac{(N+1)(p-1)}{24}-\sum_{e_i>1}\frac1{e_i}.
\]
In particular $R_{N,p}=0$ exactly when no thin supersingular value exists.
\end{corollary}
\begin{proof}
Separate the thin and non-thin terms in the mass identity of
Proposition~\ref{prop:S}; use Theorem~\ref{thm:main}(b) for the last assertion.
\end{proof}

\section{The Kaneko--Zagier step}\label{sec:kz}

The following is the level-$N$ version of \citet[\S10, (35)--(38)]{KZ98}, who credit the method to Koike.

\begin{lemma}\label{lem:dpsi}
In $\Fp((q))$, with $d/dT=(\theta T)^{-1}\theta$,
\[
\frac{d\overline\psi}{dT}=T^{p-1}-(\theta T)^{p-1},\qquad (\theta T)^{p-1}=g(\overline T),\quad
g:=\frac{(\theta T)^{p-1}}{A^2}\in\Fp(T).
\]
\end{lemma}
\begin{proof}
Over $\Z$, $\theta\psi=T^{p-1}\theta T-(\theta T)(q^p)$, and $(\theta T)(q^p)\equiv(\theta T)^p$.
Divide by $\theta T=-q^{-1}+\cdots$, which is invertible. The function $g$ is a ratio of mod-$p$
forms of weight $2(p-1)$, it is $w_N$-invariant by Lemma~\ref{lem:AW}, and it has $q$-expansion
$(\theta T)^{p-1}$ because $A\equiv1$.
\end{proof}

\begin{lemma}[local structure of $g$]\label{lem:g}
Let $t\in\Fq$ be a non-cuspidal point with effective inertia $e_t$, and put $\varepsilon_t=1$ if
$t$ is supersingular and $0$ otherwise. Then
\[
\mathrm{ord}_{T=t}\,g=\frac{(p-1)(e_t-1)-2\varepsilon_t}{e_t}.
\]
For thin $t$ and $x$ over $t$, let $a=A/(\theta T)^{(p-1)/2}$, a function on $X_0(N)$ near $x$.
Then $a$ has a simple zero at $x$,
\[
a'(t):=\frac{da}{dT}(x)=\frac{(B/12)(x)}{(\theta T)(x)^{(p+1)/2}}\ne0,
\qquad g=a'(t)^{-2}(T-t)^{-2}\bigl(1+O(T-t)\bigr).
\]
\end{lemma}
\begin{proof}
By Fact~\ref{fact:modp}(iv) and Lemma~\ref{lem:inertia}(ii), $\theta T$ vanishes to order $e_t-1$ in a
stack parameter $z$, and $A$ to order $\varepsilon_t$; also $T-t\sim z^{e_t}$. This gives the order formula.
For thin $t$, the map $X_0(N)\to X_0(N)^+$ is \'etale at $x$ and $\theta T(x)\ne0$. In
characteristic $0$, for forms $F$ and $M$ of equal weight,
$\theta(F/M)=(\partial F\cdot M-F\cdot\partial M)/M^2$, where $\partial=\theta-\frac{k}{12}E_2$ is the
Serre derivative. Apply this with $F=E_{p-1}$ and $M=(\theta T)^{(p-1)/2}$ and reduce modulo $p$; at
$x$ we have $A(x)=0$. The form $\partial E_{p-1}$ has weight $p+1$ and $q$-expansion
$\equiv\theta(1)+\frac1{12}E_2\equiv\frac1{12}E_{p+1}$, using the Kummer congruence
$2(p+1)/B_{p+1}\equiv24$ and $\sigma_p\equiv\sigma_1$. So $\partial A=B/12$ as mod-$p$ forms. Finally
$da/dT=\theta a/\theta T$. The value $B(x)\ne0$ because $B/12$ is the derivative of $A$ at its simple zero.
\end{proof}

\begin{proof}[Proof of Theorem~\ref{thm:main}(b)--(c)]
Let $t\in\Fq$. By Proposition~\ref{prop:res}, near $t$ we have
$\overline\psi=-\overline H(t^p,t)(T-t)^{-1}+O(1)$, because $T^{p^2}-T=-(T-t)(1+O(T-t))$. So the
coefficient of $(T-t)^{-2}$ in $d\overline\psi/dT$ is $\overline H(t^p,t)$. By Lemma~\ref{lem:dpsi}
it is also $-[g]_{-2}$. Hence
\[
\Res_tR_{N,p}=-\overline H(t^p,t)=[g]_{-2}.
\]
By Lemma~\ref{lem:g}, $\mathrm{ord}_tg\ge0$ unless $t$ is thin. For $p\ge5$ and $e_t\ge2$ we have
$(p-1)(e_t-1)-2\ge0$. At thin $t$, $[g]_{-2}=a'(t)^{-2}=144(\theta T)(x)^{p+1}/B(x)^2\ne0$. Every pole is
simple by Proposition~\ref{prop:res}, so $R_{N,p}=P/D_{\rm thin}$ with $P(t_i)\ne0$ at every thin
$t_i$. Since $R_{N,p}(\infty)=0$, $\deg P<\deg D_{\rm thin}$.
\end{proof}

\begin{lemma}[Robert's isogeny identity]\label{lem:Bisog}
Let $N$ be any prime $\ne p$ (not only $N\in\mathcal N$). Let $x=(E,C)$ be a supersingular point of
$X_0(N)$, $\pi:E\to E'=E/C$, and $\omega=\pi^*\omega'$. Then
$B(E',\omega')=N\,B(E,\omega)$, $A(E',\omega')=A(E,\omega)$ and
$\Delta(N\tau)(x)=N^{-12}\Delta(E',\omega')$.
\end{lemma}
\begin{proof}
The isogeny identity is \citet[Th\'eor\`eme B; \S3.3, Lemma~7]{Robert80}; see also
\citet[Lemma~2.2]{Baker98}. We verify the normalization by $q$-expansions.
For a level-one form $f$ of weight $w$, $f(N\tau)$ at $(E,C,\omega)$ equals $N^{-w}f(E',\omega')$. Check
this on $(\mathrm{Tate}(q),\mu_N)$, where $E'\cong\mathrm{Tate}(q^N)$ and
$\pi^*\omega_{\rm can}'=N\omega_{\rm can}$. The form $E_2(\tau)-NE_2(N\tau)$ is a holomorphic
weight-$2$ form on $\Gamma_0(N)$ with $p$-integral coefficients. Modulo $p$, the forms
$B(\tau)-NB(N\tau)$ and $A\cdot(E_2-NE_2(N\tau))$ have weight $p+1$ and the same $q$-expansion, so they
are equal. At a supersingular point this gives $B(E,\omega)=N\cdot N^{-(p+1)}B(E',\omega')=N^{-1}B(E',\omega')$,
because $N^{p-1}=1$. The statement about $A$ follows in the same way from $A(N\tau)=A(\tau)$.
\end{proof}

\begin{proposition}[weighted residue product]\label{prop:product}
For thin $t_i$ and $x=(E,C)$ over $t_i$,
\[
r_iW_i=\frac{B^{p-1}}{\Delta_N^{\,k}}(x)=\beta(E)D_N(x)^k
=\Lp{N}c(E)c(E'),\qquad \beta(E):=\frac{B(E,\omega)^{p-1}}{\Delta(E,\omega)^{2k}}=c(E)^2.
\]
\end{proposition}
\begin{proof}
Near $x$, Proposition~\ref{prop:S} gives $S(T)=A^{p+1}/\Delta_N^k=a^{p+1}\,Q^{k}$ with
$Q:=(\theta T)^{12}/\Delta_N$, a function that is regular and nonzero at $x$. Hence
$W_i=[S/(T-t_i)^{p+1}]_{T=t_i}=a'(t_i)^{p+1}Q(x)^k$. With $r_i=a'(t_i)^{-2}$ this gives
\[
r_iW_i=a'(t_i)^{p-1}Q(x)^k=\Bigl(\frac{B}{12}\Bigr)^{p-1}(\theta T)^{-\frac{p^2-1}2}\cdot
\frac{(\theta T)^{\frac{p^2-1}{2}}}{\Delta_N^{k}}=\frac{B^{p-1}}{\Delta_N^k}(x).
\]
The identity $\Delta_N^k=\Delta^{2k}D_N^{-k}$ gives the second equality. For the third,
Lemma~\ref{lem:Bisog} gives $\Delta_N(x)=N^{-12}\Delta(E,\omega)\Delta(E',\omega')$ with
$N^{12k}=N^{(p^2-1)/2}=1$, and $B(E)^{(p-1)/2}=N^{-(p-1)/2}B(E')^{(p-1)/2}=\Lp NB(E')^{(p-1)/2}$.
\end{proof}

\section{The supersingular factorial evaluation}\label{sec:factorial}
Put $h=(p-1)/2$ and $m=(p+1)/2$, so $m+h=p$ and
$m+h(p+1)=pm$. We prove Theorem~\ref{thm:factorial} by comparing two
expressions for an iterated theta operator applied to $F=\eta^{p+1}$.
The leading coefficient calculation is the one underlying Serre's
functional equation for Eisenstein series, as presented by
\citet[\S2.5, Th\'eor\`eme~2, p.~119]{Robert80}. The choice of $F$ makes
its quadratic twist explicit at every $p\ge5$.

\subsection{A form at full level twelve}
Work on a geometrically connected component $X(12)$ of the compactified
full-level modular curve, over $R=\Z_{(p)}[\zeta_{12}]$ localized at a
prime above $p$. This is a smooth proper curve carrying the Hodge line
bundle $\omega$. Forms of weight $w$ are sections of $\omega^{\otimes w}$;
a choice of nonzero invariant differential on an elliptic curve allows
evaluation of such a section. The cusp parameter is $q^{1/12}$.
We use the integral $q$-expansion principle on this prime-to-$p$ level
\citep[\S1.6]{Katz73}; see also \citealp{DR73}.

\begin{lemma}\label{lem:eta-level12}
The function $\eta^2$ is a holomorphic modular form of weight one with
trivial character on $\Gamma(12)$. In particular
$F=\eta^{p+1}\in M_m(\Gamma(12))$ is integral over $R$, and
\begin{equation}\label{eq:factorial-Fpowers}
F^{24}=\Delta^{p+1},\qquad F^{p-1}=\Delta^k.
\end{equation}
\end{lemma}
\begin{proof}
The eta transformation law \citep[Ch.~4]{Knopp70} gives
$\eta^2(\gamma\tau)=\nu(\gamma)^2(c\tau+d)\eta^2(\tau)$.
For $\gamma=\begin{pmatrix}a&b\\c&d\end{pmatrix}\in\Gamma(12)$ with
$c>0$, the multiplier formula for $d$ odd (as here, since $12\mid c$) is
\[
\nu(\gamma)^2=
\exp\!\left(\frac{2\pi i}{12}
\bigl((a+d)c-bd(c^2-1)+3d-3-3cd\bigr)\right)=1.
\]
Here the squared Kronecker-symbol factor is one, and the displayed
numerator is divisible by $12$. For $c<0$, apply the same formula to
$-\gamma$: its numerator is $-6$ modulo $12$, so its multiplier squared
is $-1$; the factor $-c\tau-d$ supplies the other minus sign.
For $c=0$, $\gamma$ is a translation by a multiple of $12$ and the
assertion follows from $\eta^2(\tau+12)=\eta^2(\tau)$.
The form is holomorphic at every cusp. The transformation law and the
product expansion give integral cusp expansions over $R$ (the multipliers
are twelfth roots of unity), so the $q$-expansion principle supplies the
integral section. Raising to the $m$th power gives $F$.
Both equalities in \eqref{eq:factorial-Fpowers} follow from
$\eta^{24}=\Delta$ and $24k=p^2-1$.
\end{proof}

\subsection{The iterated operator and the twist}
Let $\partial_w=\theta-\frac{w}{12}E_2$ be the Serre derivative.
It sends $p$-integral forms of weight $w$ to $p$-integral forms of
weight $w+2$. Define
\[
G_0=F,\qquad
G_{i+1}=E_{p-1}\partial_{w_i}G_i+\frac{w_i}{12}E_{p+1}G_i,
\qquad w_i=m+i(p+1).
\]
Then $G_i$ has weight $w_i$. The Kummer congruences
$E_{p-1}\equiv1$ and $E_{p+1}\equiv E_2\pmod p$ show that its
$q$-expansion is $\theta^iF\pmod p$. Moreover
\begin{equation}\label{eq:factorial-leading}
G_i=\gamma_i E_{p+1}^{\,i}F+E_{p-1}R_i,\qquad
\gamma_i=12^{-i}\prod_{j=0}^{i-1}w_j,
\end{equation}
with $R_i$ a $p$-integral modular form. This follows by induction: in the
recursion the derivative term is multiplied by $E_{p-1}$, and the
remaining term multiplies the leading coefficient by $w_i/12$.
At $i=h$, Wilson's theorem gives
\begin{equation}\label{eq:factorial-gamma}
\overline\gamma_h
=12^{-h}\prod_{j=0}^{h-1}(m+j)
=12^{-h}\frac{(p-1)!}{h!}
=-\frac{12^{-h}}{h!}\quad\text{in }\Fp.
\end{equation}

Write $\chi_{12}(n)=\leg{12}{n}$ for the Kronecker character, with
values $1,-1,0$ on $n\equiv\pm1,\pm5$, and nonunits modulo $12$,
respectively. Euler's pentagonal identity \citep{Apostol90} takes the form
\[
\eta(\tau)=\sum_{n\ge1}\chi_{12}(n)q^{n^2/24}.
\]
\begin{lemma}[quadratic twist]\label{lem:factorial-twist}
The following identity holds between mod-$p$ forms of weight $pm$ on
$\Gamma(12)$:
\begin{equation}\label{eq:factorial-twist}
\overline G_h=\Lp6\bigl(A^m\overline F-\chi_{12}(p)\overline F^{\,p}\bigr).
\end{equation}
\end{lemma}
\begin{proof}
By the product expansion and the freshman's dream,
$F\equiv\eta(\tau)\eta(p\tau)\pmod p$. Thus
\[
\overline F=\sum_{a,b\ge1}\chi_{12}(a)\chi_{12}(b)
q^{(a^2+pb^2)/24}.
\]
Applying $\theta^h$ multiplies each term by
$((a^2+pb^2)/24)^h$, which is zero modulo $p$ when $p\mid a$
and equals $\leg{24}{p}=\Lp6$ otherwise.
The sub-sum with $a=pa'$ is
\[
\chi_{12}(p)\sum_{a',b\ge1}\chi_{12}(a')\chi_{12}(b)
q^{p(b^2+pa'^2)/24}
=\chi_{12}(p)\overline F^{\,p}.
\]
Consequently the $q$-expansion of $\overline G_h$ equals
$\Lp6(\overline F-\chi_{12}(p)\overline F^{\,p})$.
The integral forms $G_h$, $E_{p-1}^mF$ and $F^p$ all have weight
$pm$. Since $A$ has $q$-expansion one, the $q$-expansion principle
proves \eqref{eq:factorial-twist} as an identity of sections.
\end{proof}

\begin{proof}[Proof of Theorem~\ref{thm:factorial}]
Let $(E,\omega_E)$ be supersingular and choose a full level-$12$
structure $\alpha$ with the fixed Weil pairing. Evaluate
\eqref{eq:factorial-leading} and \eqref{eq:factorial-twist} at
$x=(E,\alpha,\omega_E)$. Since $A(x)=0$, they give
\[
-\frac{12^{-h}}{h!}\,B(E,\omega_E)^h\overline F(x)
=-\Lp6\chi_{12}(p)\overline F(x)^p.
\]
Equation~\eqref{eq:factorial-Fpowers} shows that
$\overline F(x)\ne0$ and
$\overline F(x)^{p-1}=\Delta(E,\omega_E)^k$. Hence
\[
B(E,\omega_E)^h
=\Lp6\chi_{12}(p)12^h h!\,\Delta(E,\omega_E)^k.
\]
Euler's criterion gives $12^h=\leg{12}{p}=\chi_{12}(p)$,
so the two character factors cancel. This proves the pointwise formula
for every supersingular curve, including non-rational $j$-invariants.

For the divisibility statement, work in $\Fb[E_4,E_6]$.
The Hasse invariant $A$ is squarefree. Its factors are $E_4$ if
$j=0$ is supersingular, $E_6$ if $j=1728$ is supersingular, and
$E_4^3-j_0\Delta$ for the remaining supersingular $j_0$
\citep[\S3]{KZ98}. This is the geometric factorization over $\Fb$;
the irreducible quadratic factors in $j$ over $\Fp$ in
\citet[Theorem~1.2]{Baker98} split after extending the constants.
For $j_0\ne0,1728$, the factor
\[
E_4^3-j_0\Delta
=\left(1-\frac{j_0}{1728}\right)E_4^3
+\frac{j_0}{1728}E_6^2
\]
is irreducible: it has the form $aX^3+bY^2$ with $ab\ne0$.
The zero set of each factor away from the origin
parametrizes the corresponding curves with their nonzero differentials.
The pointwise identity therefore forces each factor to divide
$B^h-\Lp6h!\Delta^k$, so their product $A$ divides it.
Divisibility descends to $\Fp[E_4,E_6]$. Conversely, divisibility
immediately implies the pointwise identity because $A$ vanishes at every
supersingular curve. The factors $E_4$ and $E_6$ account for the elliptic
points $j=0,1728$ without extra multiplicities in the graded ring.
\end{proof}

\begin{proof}[Proof of Corollary~\ref{cor:D}]
Theorem~\ref{thm:factorial} gives $c(E)=\Lp6h!\in\Fp^\times$
independently of $E$. Squaring and using
$(h!)^2=(-1)^{h+1}=-\Lp{-1}$ proves the stated form of
Baker's congruence. This proof uses neither that congruence nor
connectivity of supersingular isogeny graphs as an input.
\end{proof}

\section{The discriminant ratio and signed residues}\label{sec:proof}
We use the algebraic unit $D_\ell$ of \eqref{eq:Dmoduli}, with
$\pi:E\to E'=E/C$ and $\omega=\pi^*\omega'$.

\begin{proof}[Proof of Theorem~\ref{thm:C}]
At a supersingular point $x=(E,C)$, Lemma~\ref{lem:Bisog} and
\eqref{eq:Dmoduli} imply
\[
\frac{c(E')}{c(E)}
=\ell^h\left(\frac{\Delta(E,\omega)}{\Delta(E',\omega')}\right)^k
=\Lp\ell\,D_\ell(x)^k,
\]
because $\ell^{12k}=\ell^{(p^2-1)/2}=1$.
By Corollary~\ref{cor:D}, the left side is one and is nonzero.
Therefore $D_\ell(x)^k=\Lp\ell$. For $\ell=N\in\mathcal N$,
$D_N=u_N^{N-1}$ and $(N-1)k=(p^2-1)/a_N$, giving the assertion for $u_N$.
\end{proof}

\begin{remark}[comparison with Ono]\label{rem:ono}
\citet[Theorem~1.2 and Remark~4]{Ono25} uses the reciprocal
$R_\ell=\Delta(\ell\tau)/\Delta(\tau)=D_\ell^{-1}$.
Since a Legendre symbol equals its inverse, Theorem~\ref{thm:C}
determines the same half-exponent sign for $R_\ell$, namely $\Lp\ell$.
For example, $\leg2{11}=-1$, as in his example.
His Theorem~1.3 states that on a model with $p^2$-Frobenius $[-p]$
and an $\Fq$-rational differential,
$\Delta(E,\omega)^{2k}=(-1)^{(p+1)/2}$.
Theorem~\ref{thm:factorial} instead relates $B(E,\omega)^h$ to
$\Delta(E,\omega)^k$ for every nonzero differential.
The ratio argument requires no descent of a twenty-fourth root of
$\Delta$ to $\Fq$.
\end{remark}

\begin{proof}[Proof of Theorem~\ref{thm:main}(e)]
Proposition~\ref{prop:product} and Corollary~\ref{cor:D} give directly
\[
r_iW_i=\Lp N\,c(E)c(E')
=\Lp N\,(\Lp6h!)^2=-\Lp{-N}.
\]
By Corollary~\ref{cor:rFp}, $r_i\in\Fp^\times$, and hence
$W_i=-\Lp{-N}r_i^{-1}\in\Fp^\times$.
\end{proof}

\begin{proof}[Proof of Theorem~\ref{thm:main}(f)]
Proposition~\ref{prop:res} gives
$\overline H(t^p,t)=-\Res_tR_{N,p}$. By part (b) this is zero
unless $t=t_i$ is thin, in which case it equals
$-r_i=\Lp{-N}W_i^{-1}$ by part (e).
\end{proof}

\begin{remark}[$\Fp$-rationality at non-rational points]
The two rationality assertions in Theorem~\ref{thm:main}(e) have different
explanations. The assertion $r_i\in\Fp$ follows from symmetry of the
Hecke polynomial and Frobenius-equivariance of residues
(Corollary~\ref{cor:rFp}). In contrast, writing
\[
W_i=\operatorname{Nm}_{\Fq/\Fp}(D_{\rm thin}'(t_i))
\prod_{e_k>1}(t_i-t_k)^{(p+1)/e_k},
\]
shows that $W_i\in\Fp$ requires the product of the thick-point factors
to lie in $\Fp$. This is not visible factor by factor; it follows here
from the signed product identity together with Corollary~\ref{cor:rFp}:
$W_i=-\Lp{-N}/r_i\in\Fp^\times$.
\end{remark}


\section{Geometric interpretation of the weights}\label{sec:geometry}
The proof of Theorem~\ref{thm:main} is algebraic. We now interpret its weights in
the integral modular curve, retaining the geometric explanation of pole cancellation.
Put $W=W(\Fb)$ and let
\[
\mathcal B=X_0(Np)^{\mathrm{DR}}_W,\qquad
\mathcal A=\mathcal B/\langle w_N\rangle.
\]
Here $\mathcal B$ denotes the coarse Deligne--Rapoport model and $w_N$ is the
partial Atkin--Lehner involution. A split node with completed local ring
$W[[s,t]]/(st-p^ev)$, $v$ a unit, has \emph{thickness} $e$.

\begin{proposition}\label{prop:geometry}
The special fibre of $\mathcal A$ has two components, each isomorphic to
$X_0(N)^+_{\Fb}$. They meet at the supersingular Fricke orbits. At the orbit of
$x=(E,C_N)$ the node has thickness
\[
e_x=|\Aut(E,C_N)/\{\pm1\}|\,
|\Stab_{\langle w_N\rangle}(x)|.
\]
Consequently $R_{N,p}$ has a pole at a supersingular value precisely when the
corresponding node has thickness one.
\end{proposition}
\begin{proof}
Deligne--Rapoport's theorem applies with prime-to-$p$ level $N$ and the stabilizer
of a line in $(\Z/N)^2$ as the auxiliary group
\citep[V, 1.14--1.20; VI, Theorem~6.9, especially (6.9.1)]{DR73}.
It gives two copies of $X_0(N)_{\Fb}$ with Frobenius gluing, and the coarse
completed local ring at $x$ is
\[
W[[s,t]]/(st-p^{m_x}),\qquad m_x=|\Aut(E,C_N)/\{\pm1\}|.
\]
The forgetful map $\mathcal B\to X_0(N)_W$ is finite: it is proper and has finite
fibres on the generic fibre and on the two special-fibre components.
Normality identifies $\mathcal B$ with the normalization in its generic function
field. The generic partial Fricke map covers $w_N$ on $X_0(N)_W$, hence extends
uniquely by normalization. On the elliptic locus it is
\[
(E,C_N,C_p)\longmapsto(E/C_N,E[N]/C_N,q_N(C_p)).
\]
Since $q_N$ is prime to $p$, it preserves the Frobenius and Verschiebung components
and both branches at a node.

To apply Liu's quotient result, note that $\mathcal B$ is a normal,
proper flat fibred surface over the DVR $W$, with smooth generic fibre
and nodal special fibre. At a fixed geometric node its stabilizer has
order two, prime to $p$, and acts trivially on its residue field $\Fb$.
As just shown it preserves both branches. It acts nontrivially on either
coarse branch. Indeed its action on the smooth curve $X_0(N)$ is a nontrivial
tame involution, whose derivative at a fixed point is $-1$.
Taking the branch-preserving quotient therefore multiplies the thickness by
this order, by \citet[Proposition~10.3.48(b)--(c), pp.~526--528]{Liu02}.
At a non-fixed node the inertia is trivial and the thickness is unchanged.
Invariants commute with reduction because $2$ is invertible in $W$.
The quotient components are consequently $X_0(N)^+_{\Fb}$, and no node is
smoothed, since no stabilizer exchanges branches. This proves the thickness formula.
The last assertion is Theorem~\ref{thm:main}(b).
\end{proof}

\begin{remark}[level two]
For $N=2$, the Fricke fixed points are $u=\pm64$, with $T=-104,152$.
At $u=-64$ the reduced automorphism group has order two and the Fricke stabilizer
has order two, giving thickness four. At $u=64$ their orders are one and two,
giving thickness two. All other supersingular values have thickness one.
The two exceptional values are supersingular precisely when
$\left(\frac{-1}{p}\right)=-1$ and $\left(\frac{-2}{p}\right)=-1$, respectively.
Thus a thick value is absent from the pole divisor but remains in the
weighted residue product. The same distinction applies at each of the five levels.
\end{remark}

\section{Finite certificates and computational evidence}\label{sec:comp}
The accompanying exact-arithmetic computations provide checks independent of the
signed product formula. They are not used in the proof. We describe the certification
bound so that the strength of each computation is explicit.

\begin{proposition}[finite certificate]\label{prop:certificate}
Let $L=p^2$, let $D\in\Fp[T]$ be a monic divisor of $T^L-T$, of degree $d$, and
let $P\in\Fp[T]$ have degree less than $d$. If
\[
D(\overline T)\,\overline{T|U_p}-P(\overline T)=O(q^{L-d+1}),
\]
then $R_{N,p}=P/D$. For $D=1$, take $P=0$.
\end{proposition}
\begin{proof}
Proposition~\ref{prop:res} implies that
$R_{N,p}-P/D=C/(T^L-T)$ with $\deg C<L$.
If $C\ne0$, its $q$-order is $L-\deg C$, since $T=q^{-1}+O(q)$.
Multiplication by $D(T)$ gives order $L-\deg C-d\le L-d$, contradicting the
hypothesis. Thus $C=0$.
\end{proof}

The first computation reconstructs $R_{N,p}$ from its $q$-expansion, using
$D=D_{\rm thin}$ built independently from supersingular data.
The coefficients of the numerator are recovered without the residue formula,
and Proposition~\ref{prop:certificate} then proves the reconstructed identity.
The reduced denominator was also recovered without a candidate for $p\le37$.
The Hecke polynomial was computed modulo $p^2$ for $p\le31$, checking symmetry,
the Kronecker congruence, and the residue criterion at every element of $\Fq$.

\begin{table}[ht]\centering
\begin{tabular}{rrrrr}\toprule
$N$ & certified pairs & thin values & non-rational thin values & range\\\midrule
2 & 45 & 530 & 310 & $5\le p\le211$\\
3 & 45 & 716 & 456 & $5\le p\le211$\\
5 & 44 & 1059 & 740 & $5\le p\le211$, $p\ne5$\\
7 & 44 & 1431 & 1116 & $5\le p\le211$, $p\ne7$\\
13 & 44 & 2510 & 2074 & $5\le p\le211$, $p\ne13$\\\bottomrule
\end{tabular}
\caption{Certified reconstructions. Thin values are distinct geometric values of $T_N$.}
\end{table}

All 222 admissible pairs with $p\le211$ passed the certificates, giving 6246 thin
values, of which 4696 are non-rational. Four pairs have no thin values; the ratio
set is empty in these cases. In every other pair it is a singleton equal to
$-\left(\frac{-N}{p}\right)$. A second calculation, using the double-pole formula
of Lemma~\ref{lem:g}, checks a further 960 pairs for $223\le p\le1499$.
That calculation depends on the derivative lemma and is not a separate
reconstruction from $T|U_p$.

The eta/Frobenius identity was also tested by a different route: normalized
V\'elu quotients and their discriminants, without eta functions or Hauptmoduln.
For the 15 primes $\ell\le47$ and all admissible $5\le p\le199$, the calculation
covers 647 pairs and 127541 cyclic-kernel checks on chosen representatives of
supersingular elliptic curves. These are not counts of distinct coarse points:
automorphisms may identify kernels over $j=0$ and $1728$.
Every kernel gave the asserted discriminant-ratio sign.

For Theorem~\ref{thm:factorial}, the supplementary computations cover all
166 primes $5\le p\le997$. Exact polynomial remainders and evaluations
on 6430 supersingular $j$-values, of which 4734 lie outside $\Fp$, agree
with the asserted constant. The curve evaluations use three randomly
rescaled models per value; additional routes use Legendre and Weierstrass
coefficients. Replacing the constant by its negative fails at every prime
in this range. The two theta identities used in the proof were also
checked as polynomial identities in $\Fp[E_4,E_6]$ at
$p=23,47,71,167$, where $F=\eta^{p+1}$ has level one.
These finite checks are independent evidence, not premises of the proof.

\paragraph{Data and software.}
The supplementary packages \texttt{Signed\_Weighted\_Residue\_v1\_0} and
\texttt{Eta\_Frobenius\_Lemma\_Audit\_v1\_0} contain Python code, per-pair JSON
results and reproduction instructions (\texttt{python-flint} 0.9.0 and
\texttt{mpmath}). The package \path{Factorial_Conjecture_Result_v1_0}
contains the factorial-evaluation code and its recorded results under
\texttt{new\_work/}, with reproduction commands in \texttt{run\_all.sh}
(\texttt{cypari2} 2.2.4, PARI/GP 2.17.2 and \texttt{mpmath} 1.3.0).
The reported ranges are those documented by the exact-arithmetic results
in these supplementary packages.
\if\relax\detokenize\expandafter{\DataRepositoryURL}\relax
The code, per-pair data and reproduction instructions are supplied as
supplementary files with the manuscript.
\else
The code, data and reproduction instructions are publicly available at
\expandafter\url\expandafter{\DataRepositoryURL}.
\fi

\section{Further directions}
The argument separates three inputs: a genus-zero coordinate and its
Kronecker congruence, the Hasse divisor and Kodaira--Spencer derivative,
and the factorial evaluation of Theorem~\ref{thm:factorial}.
The first two give the pole criterion in the chosen coordinate; the
third selects the residue sign and also yields the discriminant-ratio
identity of Theorem~\ref{thm:C} for every prime isogeny degree different
from $p$. Extending the residue theorem to other Fricke curves requires
an integral normalized coordinate and an explicit treatment of the cusp
divisor; it is not asserted here.

The proof of Theorem~\ref{thm:factorial} suggests seeking other modular
forms whose quadratic twists admit an equally explicit expression.
Such an extension would require a counterpart of the twist identity,
in addition to the standard theta-iteration argument. The genus-zero
moonshine identification supplies examples of the modular functions
to which the residue theorem applies; the result itself is a statement
about their reduction and supersingular divisors.

\paragraph{Acknowledgement of assistance.}
AI systems assisted in developing arguments, checking calculations and
preparing this manuscript, including the factorial-evaluation argument
and its exact-arithmetic verification code. The source computations
and proof checks are documented in the supplementary packages.
This assistance does not constitute independent specialist refereeing.

\begin{thebibliography}{99}\small
\bibitem[Apostol(1990)]{Apostol90} Apostol, T.~M. (1990). \emph{Modular Functions and Dirichlet Series in Number Theory}, 2nd ed. Graduate Texts in Mathematics 41. Springer.

\bibitem[Baker(1998)]{Baker98} Baker, A. (1998). A supersingular congruence for modular forms.
\emph{Acta Arithmetica}, 86(1), 91--100. \url{https://doi.org/10.4064/aa-86-1-91-100}.

\bibitem[Behrens(2009)]{Behrens09} Behrens, M. (2009). Congruences between modular forms given by the divided $\beta$ family in homotopy theory. \emph{Geometry \& Topology}, 13, 319--357. \url{https://doi.org/10.2140/gt.2009.13.319}.

\bibitem[Betina and Lecouturier(2018)]{BL18} Betina, A., \& Lecouturier, E. (2018). Congruence formulas for
Legendre modular polynomials. \emph{Journal of Number Theory}, 188, 71--87.

\bibitem[Bringmann et~al.(2026)]{BGKMR26} Bringmann, K., Guerzhoy, P.,
Kane, B., Mertens, M., \& Rolen, L. (2026). On $U_p$-congruences for
meromorphic modular forms via supersingularity. \emph{arXiv:2606.14020v1}.
\url{https://arxiv.org/abs/2606.14020v1}.

\bibitem[Chen et~al.(2019)]{CMT19} Chen, R.~C., Marks, S., \& Tyler, M. (2019).
$p$-Adic properties of Hauptmoduln with applications to moonshine.
\emph{SIGMA}, 15, 033, 35 pp. \url{https://doi.org/10.3842/SIGMA.2019.033}.


\bibitem[Conway and Norton(1979)]{CN79} Conway, J.~H., \& Norton, S.~P. (1979). Monstrous moonshine.
\emph{Bulletin of the London Mathematical Society}, 11, 308--339.

\bibitem[Deligne and Rapoport(1973)]{DR73} Deligne, P., \& Rapoport, M. (1973). Les sch\'emas de modules de
courbes elliptiques. In \emph{Modular Functions of One Variable II} (LNM 349, pp.~143--316). Springer.

\bibitem[de~Shalit(1994)]{dS94} de~Shalit, E. (1994). Kronecker's polynomial, supersingular elliptic curves,
and $p$-adic periods of modular curves. \emph{Contemporary Mathematics}, 165, 135--148.

\bibitem[Deuring(1941)]{Deuring41} Deuring, M. (1941). Die Typen der Multiplikatorenringe elliptischer Funktionenk\"orper. \emph{Abhandlungen aus dem Mathematischen Seminar der Hansischen Universit\"at}, 14, 197--272.

\bibitem[Duncan and Swisher(2026)]{DS26} Duncan, J.~F.~R., \& Swisher, H. (2026).
Modular functions and the monstrous exponents. \emph{arXiv:2602.09135v1}, 9~February~2026. \url{https://arxiv.org/abs/2602.09135v1}.

\bibitem[Gross(1990)]{Gross90} Gross, B.~H. (1990). A tameness criterion for Galois representations associated to modular forms (mod $p$). \emph{Duke Mathematical Journal}, 61(2), 445--517.

\bibitem[Jorgenson et~al.(2019)]{JST19} Jorgenson, J., Smajlovi\'c, L., \& Then, H. (2019).
The Hauptmodul at elliptic points of certain arithmetic groups.
\emph{Journal of Number Theory}, 204, 661--682.
\url{https://doi.org/10.1016/j.jnt.2019.03.021}.
The elliptic-value tables used here are also available in
\emph{arXiv:1602.07426v2}, \S4.2, Table~1 and Appendix~C.
\url{https://arxiv.org/abs/1602.07426v2}.

\bibitem[Kaneko and Zagier(1998)]{KZ98} Kaneko, M., \& Zagier, D. (1998). Supersingular $j$-invariants,
hypergeometric series, and Atkin's orthogonal polynomials. In \emph{Computational Perspectives on Number
Theory} (AMS/IP Studies in Advanced Mathematics 7, pp.~97--126).

\bibitem[Katz(1973)]{Katz73} Katz, N.~M. (1973). $p$-adic properties of modular schemes and modular forms.
In \emph{Modular Functions of One Variable III} (LNM 350, pp.~69--190). Springer.


\bibitem[Katz and Mazur(1985)]{KM85} Katz, N.~M., \& Mazur, B. (1985). \emph{Arithmetic Moduli of Elliptic
Curves} (Annals of Mathematics Studies 108). Princeton University Press.

\bibitem[Knopp(1970)]{Knopp70} Knopp, M.~I. (1970). \emph{Modular Functions in Analytic Number Theory}. Markham.

\bibitem[Koike(1973)]{Koike73} Koike, M. (1973). Congruences between modular forms
and functions and applications to the conjecture of Atkin.
\emph{Journal of the Faculty of Science, University of Tokyo, Section IA, Mathematics},
20, 129--169.

\bibitem[Landweber(1988)]{Landweber88} Landweber, P.~S. (1988). Supersingular elliptic curves and congruences for Legendre polynomials. In \emph{Elliptic Curves and Modular Forms in Algebraic Topology}, Lecture Notes in Mathematics 1326, pp.~69--93. Springer.

\bibitem[Larson(2020)]{Larson20} Larson, D.~M. (2020). Explicit modular forms from the divided beta family. \emph{Topology and its Applications}, 273, 107094. \url{https://doi.org/10.1016/j.topol.2020.107094}.

\bibitem[Liu(2002)]{Liu02} Liu, Q. (2002). \emph{Algebraic Geometry and Arithmetic Curves}. Oxford Graduate Texts in Mathematics 6. Oxford University Press.


\bibitem[Morton(2023)]{Morton23} Morton, P. (2023). Supersingular conjectures for the
Fricke group. \emph{International Journal of Number Theory}, 19(4), 803--828.
\url{https://doi.org/10.1142/S1793042123500409}; preprint
\url{https://arxiv.org/abs/2105.14532}.

\bibitem[Nakaya(2019)]{Nakaya19} Nakaya, T. (2019). The number of linear factors of
supersingular polynomials and sporadic simple groups. \emph{Journal of Number Theory},
204, 471--496. \url{https://doi.org/10.1016/j.jnt.2019.04.015}.


\bibitem[Ono(2026)]{Ono25} Ono, K. (2026). Modular forms for chromatic homotopy: supersingular congruences. \emph{arXiv:2509.16175v2}, 3 July 2026. \url{https://arxiv.org/abs/2509.16175v2}.


\bibitem[Robert(1980)]{Robert80} Robert, G. (1980). Congruences entre s\'eries d'Eisenstein, dans le cas supersingulier. \emph{Inventiones Mathematicae}, 61, 103--158.

\bibitem[Sakai(2011)]{Sakai11} Sakai, Y. (2011). The Atkin orthogonal polynomials for
the low-level Fricke groups and their application.
\emph{International Journal of Number Theory}, 7(6), 1637--1661.
\url{https://doi.org/10.1142/S1793042111004460}.

\bibitem[Sakai(2014)]{Sakai14} Sakai, Y. (2014). The Atkin orthogonal polynomials for
the Fricke groups of levels $5$ and $7$.
\emph{International Journal of Number Theory}, 10(8), 2243--2255.
\url{https://doi.org/10.1142/S1793042114500766}.

\bibitem[Waterhouse(1969)]{Waterhouse69} Waterhouse, W.~C. (1969). Abelian varieties over finite fields. \emph{Annales Scientifiques de l'\'Ecole Normale Sup\'erieure (4)}, 2, 521--560.
\end{thebibliography}
\end{document}
